\section{\sname: Coordinated Escape and Barrier Safety in One Nav2 Controller}\label{sec:method}

\noindent\textbf{Problem Setup.} Consider a differential-drive robot with pose $\mathbf{x}_t=(x_t,y_t,\theta_t)$, position $\mathbf{p}_t=(x_t,y_t)$ and control $\mathbf{u}_t=(v_t,\omega_t)\in\mathcal{U}$, where $\mathcal{U}$ is the box $v\in[v_{\min},v_{\max}]$, $|\omega|\le\omega_{\max}$, following the unicycle model $\dot x=v\cos\theta$, $\dot y=v\sin\theta$, $\dot\theta=\omega$. A Nav2 global planner supplies a path $\sigma=(\sigma_1,\dots,\sigma_M)$. Around the robot are static structure, seen only through the local costmap, and a set $\mathcal{O}_t$ of dynamic obstacles $o$ with position $\mathbf{p}_o$, radius $R_o$ and estimated velocity $\mathbf{v}_o$. Let $r$ be the robot's inscribed radius and $m$ a safety margin. We seek a controller that (G1) keeps making progress along $\sigma$ instead of stalling in a local minimum, (G2) keeps the robot clear of the obstacles in $\mathcal{O}_t$ through a barrier whose forward invariance holds under stated conditions, and (G3) runs as a Nav2 controller in real time on a CPU. Formally, each control cycle approximately solves
\begin{equation}\label{eq:objective}
\begin{aligned}
\mathbf{u}^\star_{t:t+T-1}=\;&\arg\min_{\mathbf{u}_{t:t+T-1}\in\,\mathcal{U}^{T}}\;
\mathbb{E}\Big[\sum_{\tau=t}^{t+T-1} c\big(\mathbf{x}_\tau,\mathbf{u}_\tau;\sigma\big)
+ e_t\,c_{\mathrm{esc}}\big(\mathbf{x}_{t:t+T}\big)\Big]\\
&\text{s.t.}\quad h_o\ \text{satisfies the CBF condition \eqref{eq:cbfcond} at }(\mathbf{x}_t,\mathbf{u}_t)\quad\forall o\in\mathcal{O}_t ,
\end{aligned}
\end{equation}
where $c$ is the stock Nav2 critic cost, $e_t\in\{0,1\}$ is the entrapment flag, $c_{\mathrm{esc}}$ are the escape costs, and $h_o$ is a control barrier function for obstacle $o$. The tension is that G1 may demand manoeuvres that approach obstacles to round a trap, whereas a fixed safety filter for G2 forbids exactly those manoeuvres. \sname{} splits \eqref{eq:objective}: MPPI minimizes the expected cost, with escape costs switched in by $e_t$ (\cref{sec:sample,sec:detect,sec:escape}); a per-cycle QP enforces the barrier on the first control (\cref{sec:track,sec:cbf}); and a coordinator couples the two through the barrier gain (\cref{sec:coord}). \Cref{alg:semppi} summarizes one cycle.

\subsection{Sampling Nominal Commands: Reusing the Nav2 MPPI Optimizer}\label{sec:sample}

\sname{} subclasses the stock \code{MPPIController} and calls its optimizer unchanged. Given the previous control sequence $\bar{\mathbf{u}}_{0:T-1}$, the optimizer draws $K$ perturbations $\boldsymbol{\epsilon}^k_j\sim\mathcal{N}(\mathbf{0},\Sigma)$, rolls out $\mathbf{u}^k_j=\mathrm{clip}_{\mathcal{U}}(\bar{\mathbf{u}}_j+\boldsymbol{\epsilon}^k_j)$ through the motion model, and scores each rollout $\mathbf{x}^k_{0:T}$ with the listed critics $C_i$, among which the \code{EscapeCritic} (\cref{sec:escape}):
\begin{equation}\label{eq:mppi}
\begin{aligned}
S^k&=\sum_i C_i\big(\mathbf{x}^k_{0:T}\big)+e_t\Big(C_{\mathrm{apf}}\big(\mathbf{x}^k_{0:T}\big)+C_{\mathrm{gap}}\big(\mathbf{x}^k_{0:T}\big)\Big),\\
w^k&=\frac{\exp\!\big(-(S^k-\min_{k'}S^{k'})/\lambda\big)}{\sum_{k''}\exp\!\big(-(S^{k''}-\min_{k'}S^{k'})/\lambda\big)},\qquad
\bar{\mathbf{u}}_j\leftarrow\sum_k w^k\,\mathbf{u}^k_j .
\end{aligned}
\end{equation}
The first element is the nominal command $\mathbf{u}^{\mathrm{nom}}_t=(\bar v_0,\bar\omega_0)$ \citep{mppi_orig}. (Nav2 additionally regularizes the control effort through a small weight $\gamma$; we omit it for brevity.) With the released defaults, $T=56$, $\Delta t=0.05$\,s, $K=1000$, $\lambda=0.3$ and $\Sigma^{1/2}=\mathrm{diag}(0.2,0.4)$ (\cref{tab:params}).

\subsection{Detecting Entrapment: A Single Monotone Progress Signal}\label{sec:detect}

The nearest-path index $k_t=\arg\min_i\|\mathbf{p}_t-\sigma_i\|$ is non-monotone: it jumps back when the robot backs off a wall. We therefore track the furthest index ever reached and count the cycles in which it fails to advance,
\begin{equation}\label{eq:detect}
\begin{gathered}
k^{\max}_t=\max\big(k^{\max}_{t-1},k_t\big),\qquad
s_t=\begin{cases}0 & k^{\max}_t>k^{\max}_{t-1},\\ s_{t-1}+1 & \text{otherwise},\end{cases}\\
e_t=\begin{cases}0 & k^{\max}_t>k^{\max}_{t-1}\ \text{or}\ \|\mathbf{p}_t-\sigma_M\|\le\epsilon_{\mathrm{goal}},\\ 1 & s_t\ge W,\\ e_{t-1} & \text{otherwise}.\end{cases}
\end{gathered}
\end{equation}
Entrapment thus latches after $W$ stalled cycles ($W=30$ by default) and clears on the first real progress, so stop-and-go motion during an escape cannot leave the flag stuck; within $\epsilon_{\mathrm{goal}}$ of the path end the detector is reset, so a robot finishing at the goal does not trigger a false escape. Within \sname, the controller's detector is the single entrapment source: the critic and the coordinator read the same atomic flag through a registry keyed by the controller name, so escape and safety always agree on whether the robot is trapped. (When the critic runs under a stock MPPI controller, it falls back to its own detector on MPPI's furthest-reached index.)

\subsection{Escaping Local Minima: Gated Repulsion and Gap Attraction}\label{sec:escape}

Only while $e_t=1$, the \code{EscapeCritic} adds two terms to every rollout's cost in \eqref{eq:mppi}; in free space the controller is therefore exactly stock MPPI (detect-and-switch rather than always-on, as in DRPA-MPPI \citep{drpa_mppi}). The first is a distance-field artificial potential of the classical repulsive form \citep{khatib_apf}. Let $d(\mathbf{q})$ be the 8-connected Dijkstra distance from costmap cell $\mathbf{q}$ to the nearest obstacle cell, where unknown cells are \emph{not} obstacles so that unexplored space does not repel the escape. With $\mathcal{J}^k$ the on-map poses of rollout $k$ and $d_{\min}$ half a cell,
\begin{equation}\label{eq:apf}
U(d)=\begin{cases}\tfrac12\,\eta\,\Big(\dfrac{1}{\max(d,d_{\min})}-\dfrac{1}{d_0}\Big)^{2} & d<d_0,\\[4pt] 0 & d\ge d_0,\end{cases}\qquad
C_{\mathrm{apf}}\big(\mathbf{x}^k_{0:T}\big)=\Big(\frac{1}{|\mathcal{J}^k|}\sum_{j\in\mathcal{J}^k}U\big(d(\mathbf{p}^k_j)\big)\Big)^{p},
\end{equation}
which pushes rollouts off the trapping surface. Repulsion alone does not route the robot \emph{around} a non-convex trap, so the second term supplies a temporary heading toward an opening (a bearing cost; the goal itself is unchanged), a lightweight form of the free-space subgoal and gap guidance of \citet{gp_mppi} and \citet{potential_gap}. The critic casts $N_r$ rays at bearings $\psi_i=-\pi+2\pi i/N_r$ from the robot, measures each ray's free length $\ell(\psi_i)\le r_{\max}$ on the costmap, and selects the open bearing closest to the goal bearing $\psi_{\mathrm{goal}}$; rollouts whose endpoint heads toward it are rewarded:
\begin{equation}\label{eq:gap}
\psi^\star=\arg\min_{\psi_i:\;\ell(\psi_i)\ge r_{\min}}\big|\mathrm{wrap}(\psi_i-\psi_{\mathrm{goal}})\big|,\qquad
C_{\mathrm{gap}}\big(\mathbf{x}^k_{0:T}\big)=w_g\Big(1-\cos\mathrm{wrap}\big(\angle(\mathbf{p}^k_T-\mathbf{p}_t)-\psi^\star\big)\Big)^{p_g}.
\end{equation}
The defaults are $d_0=0.6$\,m, $\eta=1$, $N_r=36$, $r_{\max}=2$\,m, $r_{\min}=0.6$\,m and $w_g=4$ (\cref{tab:params}). In the 2D mirror of \cref{sec:experiments}, which uses the selected bearing as a temporary subgoal rather than through $C_{\mathrm{gap}}$ (\cref{app:mirror}), disabling the gap search removes every escape (\cref{sec:ablation}); the $C_{\mathrm{gap}}$ term of the C++ critic has not been ablated on its own.

\subsection{Tracking Dynamic Obstacles: Costmap Clusters with Gated Association}\label{sec:track}

The tracker clusters local-costmap cells at or above the inscribed-inflation threshold, summarizing each cluster $c$ by a centroid $\mathbf{c}$ and a radius. Each current cluster is associated with at most one previous track within a gate of $0.6$\,m, each track being consumable once, so a split or merge cannot manufacture a phantom velocity. The velocity is a clamped finite difference over the track history,
\begin{equation}\label{eq:track}
\mathbf{v}_o=\mathrm{clamp}_{v_{\mathrm{track}}}\!\Big(\frac{\mathbf{c}_{t}-\mathbf{c}_{t'}}{t-t'}\Big),\qquad
\mathcal{O}_t=\big\{o:\ \|\mathbf{v}_o\|\ge v_{\mathrm{dyn}},\ R_o\le R_{\max}\big\},
\end{equation}
with $t'$ the previous observation of the track and $v_{\mathrm{track}}=2$\,m/s. Only obstacles that are moving and small enough to be a movable body enter $\mathcal{O}_t$, and hence the coordinator and the CBF. Static structure is deliberately excluded and left to the MPPI obstacle critic and costmap inflation; \cref{sec:live} shows why this scoping is essential on a real costmap. The constant-velocity model is a replaceable baseline (\cref{sec:discussion}).

\subsection{Filtering Commands: A Look-Ahead-Point CBF Quadratic Program}\label{sec:cbf}

A distance barrier on the unicycle's centre has relative degree two in $\omega$. We therefore enforce safety on the look-ahead point $\mathbf{P}_t=\mathbf{p}_t+L(\cos\theta_t,\sin\theta_t)$, which $(v,\omega)$ actuates fully:
\begin{equation}\label{eq:lookahead}
\dot{\mathbf{P}}_t=G(\theta_t)\,\mathbf{u}_t,\qquad
G(\theta)=\begin{bmatrix}\cos\theta & -L\sin\theta\\ \sin\theta & L\cos\theta\end{bmatrix}.
\end{equation}
For each $o\in\mathcal{O}_t$, with $\mathbf{d}_o=\mathbf{P}_t-\mathbf{p}_o$ and effective radius $R^{\mathrm{eff}}_o=r+R_o+m$, the barrier and its derivative are
\begin{equation}\label{eq:barrier}
h_o=\|\mathbf{d}_o\|^2-\big(R^{\mathrm{eff}}_o\big)^2,\qquad
\dot h_o=2\,\mathbf{d}_o^{\top}\big(G(\theta_t)\,\mathbf{u}_t-\mathbf{v}_o\big),
\end{equation}
and the continuous-time CBF condition \citep{ames_cbf_tac,ames_cbf_ecc}
\begin{equation}\label{eq:cbfcond}
\dot h_o+\alpha_t\,h_o\ \ge\ 0
\end{equation}
is linear in $\mathbf{u}_t$. Once per control cycle we solve, with OSQP \citep{osqp},
\begin{equation}\label{eq:qp}
\min_{v,\,\omega,\,\delta\ge0}\;w_v\big(v-v^{\mathrm{nom}}_t\big)^2+w_\omega\big(\omega-\omega^{\mathrm{nom}}_t\big)^2+\rho\,\delta^2
\quad\text{s.t.}\quad
2\,\mathbf{d}_o^{\top}G(\theta_t)\,\mathbf{u}+\delta\ \ge\ -\alpha_t\,h_o+2\,\mathbf{d}_o^{\top}\mathbf{v}_o\ \ \forall o,\qquad \mathbf{u}\in\mathcal{U},
\end{equation}
where the nominal command is first clamped to $\mathcal{U}$ and a single shared slack $\delta$ keeps the QP feasible. Obstacles are pruned to the $N_{\mathrm{obs}}$ nearest by \emph{clearance} (surface gap) rather than centre distance, so a large obstacle is never dropped in favour of a smaller one whose centre is closer. If the QP can stay safe only by relaxing the barrier ($\delta>\epsilon_\delta=10^{-3}$) or fails, the controller brakes, setting $v=0$ while keeping the QP's turn (or, on solver failure, the clamped nominal turn): it stops rather than drive into an imminent collision.

Because \eqref{eq:cbfcond} is imposed only at sampling instants, strict between-sample invariance would require tightening $R^{\mathrm{eff}}_o$ by a term of order $\dot h_{\max}\Delta t$. In practice the margin $m$ absorbs this gap, and we report the empirical slack usage instead of claiming continuous-time invariance between samples.

\subsection{Coordinating Escape and Safety: Gain Scheduling under a Conditional Invariance Lemma}\label{sec:coord}

The coordinator resolves the barrier gain from the entrapment flag and the minimum time-to-collision under the current command,
\begin{equation}\label{eq:alpha}
\alpha_t=\begin{cases}\alpha_{\mathrm{escape}} & e_t=1\ \text{and}\ \mathrm{TTC}_t\ge\tau,\\ \alpha_{\mathrm{base}} & \text{otherwise},\end{cases}\qquad
\mathrm{TTC}_t=\min_{o\in\mathcal{O}_t:\,\kappa_o>0}\frac{\|\mathbf{r}_o\|-(r+R_o)}{\kappa_o},\quad
\kappa_o=-\frac{\mathbf{r}_o^{\top}(\mathbf{v}_o-v^{\mathrm{nom}}_t\mathbf{h}_t)}{\|\mathbf{r}_o\|},
\end{equation}
where $\mathbf{r}_o=\mathbf{p}_o-\mathbf{p}_t$, $\mathbf{h}_t=(\cos\theta_t,\sin\theta_t)$, $\kappa_o$ is the closing speed, $\mathrm{TTC}_t=0$ on contact and $+\infty$ if nothing approaches, and $\alpha_{\mathrm{escape}}>\alpha_{\mathrm{base}}>0$ (defaults $6$ and $2$, $\tau=1.5$\,s).\footnote{The released controller also holds $\alpha_t$ at $\alpha_{\mathrm{base}}$ while an online conformal bound on obstacle-prediction error exceeds a trust threshold, and that bound inflates $R^{\mathrm{eff}}_o$. This is the follow-up conformal margin of \cref{sec:related}, not claimed here; an implementation of \eqref{eq:alpha} alone reproduces the benchmark configurations of \cref{sec:experiments}, not the released defaults.} Raising $\alpha_t$ relaxes the lower bound $\dot h_o\ge-\alpha_t h_o$, letting the robot approach an obstacle faster as an escape requires, while the TTC override snaps the gain back to $\alpha_{\mathrm{base}}$ when a dynamic obstacle is imminent. The next result states that, under explicit conditions, this schedule does not by itself give up forward invariance. It is a known property of CBFs with time-varying class-$\mathcal{K}$ terms and follows the same argument as the relative-degree-one case of the invariance result for adaptive CBFs with time-varying penalty functions \citep[Thm.~3]{ada_cbf}, which is stated there for Lipschitz-continuous controllers, whereas our switched gain yields a discontinuous command; related time-varying and optimized decay rates appear in \citet{optdecay_cbf} and \citet{rt_cbf}. We restate it with the standard integrating-factor proof for completeness and claim no new mathematics.

\begin{assumption}[Slack-free barrier between samples]\label{as:intersample}
For a fixed obstacle $j$ and all $t\ge0$ in the interval considered, the QP \eqref{eq:qp} is feasible with $\delta(t)=0$, and the continuous-time inequality $\dot h_j(t)\ge-\alpha(t)\,h_j(t)$ holds for almost every $t$.
\end{assumption}

The second clause is an assumption, not a consequence: feasibility at discrete instants does not by itself imply the inequality between samples (the sampled-data caveat of \cref{sec:cbf}).

\begin{proposition}[Invariance under a time-varying gain; after \citealp{ada_cbf}]\label{prop:invariance}
Let $\alpha:[0,\infty)\to[\alpha_{\min},\alpha_{\max}]\subset(0,\infty)$ be any measurable, positive, bounded gain schedule, in particular the piecewise-constant switching \eqref{eq:alpha} with its TTC override. Under \cref{as:intersample}, if $h_j(0)\ge0$ then $h_j(t)\ge0$ for all $t\ge0$: the look-ahead point does not enter obstacle $j$'s inflated disc, as modelled.
\end{proposition}

\begin{proof}
Define $\varphi(t)=\exp\big(\int_0^t\alpha(s)\,ds\big)$, which is finite and strictly positive because $0<\alpha\le\alpha_{\max}$. Almost everywhere,
\begin{equation}\label{eq:proof}
\frac{d}{dt}\big[\varphi(t)\,h_j(t)\big]=\varphi(t)\big[\dot h_j(t)+\alpha(t)\,h_j(t)\big]\ \ge\ 0,
\end{equation}
so $\varphi h_j$ is non-decreasing and $\varphi(t)h_j(t)\ge\varphi(0)h_j(0)=h_j(0)\ge0$. Since $\varphi(t)>0$, $h_j(t)\ge0$ for all $t\ge0$.
\end{proof}

Positivity and boundedness are stronger than the argument needs: a nonnegative, locally integrable $\alpha$ already makes $\varphi$ finite and positive. We keep the stronger hypothesis because it matches the schedule \eqref{eq:alpha}.

\begin{corollary}\label{cor:gain}
Under \cref{as:intersample}, forward invariance of $\mathcal{C}_j=\{h_j\ge0\}$ holds for every positive bounded schedule, including the switch to $\alpha_{\mathrm{escape}}$ and the TTC override back to $\alpha_{\mathrm{base}}$. If \cref{as:intersample} holds for every row, $\mathcal{C}=\bigcap_j\mathcal{C}_j$ is invariant; the argument is per obstacle and gain-agnostic, so sharing one $\alpha_t$ across rows does not affect it.
\end{corollary}

Under \cref{as:intersample}, raising $\alpha$ therefore does not permit $h_j<0$; it only weakens the lower bound on $\dot h_j$, admitting bolder detours inside the same set. This is the design rationale for coordinating rather than stacking the two layers: in this idealized setting, the gain trades caution for manoeuvrability without giving up invariance. The guarantee is narrow. It concerns the look-ahead point, not the robot body; it covers only the tracked dynamic obstacles under the constant-velocity model, not static structure; it applies to an obstacle only while that obstacle stays in the tracked, unpruned set $\mathcal{O}_t$ (an obstacle that slows below $v_{\mathrm{dyn}}$, merges into a cluster larger than $R_{\max}$, or is pruned beyond the $N_{\mathrm{obs}}$ smallest clearances leaves the QP); it assumes the continuous-time inequality between samples; and it requires $\delta=0$, so it does not cover any cycle in which the QP uses slack. Because the slack is shared, one imminent obstacle can momentarily loosen the margin on every other obstacle in the same solve, so we report slack usage empirically: in the benchmark of \cref{sec:experiments}, the QP never uses slack on the mover-free families, and uses nonzero slack in $26$--$34$ of $50$ narrow-dynamic trials (maximum $0.47$) and in $47$--$48$ of $50$ dynamic-family trials (maximum $0.015$) per CBF-carrying configuration (\cref{tab:slack}). In cycles with $\delta>\epsilon_\delta$ the brake of \cref{sec:cbf} replaces the QP command.

\begin{algorithm}[t]
\caption{One \sname{} control cycle (\texttt{computeVelocityCommands})}\label{alg:semppi}
\DontPrintSemicolon
\KwIn{pose $\mathbf{x}_t$, global path $\sigma$, local costmap, previous state $(k^{\max}_{t-1},s_{t-1},e_{t-1})$}
\KwOut{safe command $\mathbf{u}^{\mathrm{safe}}_t$}
$\mathbf{u}^{\mathrm{nom}}_t\gets$ stock MPPI optimizer, scoring rollouts with \eqref{eq:mppi}; the \code{EscapeCritic} adds $C_{\mathrm{apf}}+C_{\mathrm{gap}}$ iff the shared flag $e_{t-1}=1$\;
$k_t\gets\arg\min_i\|\mathbf{p}_t-\sigma_i\|$;\quad update $k^{\max}_t,s_t,e_t$ by \eqref{eq:detect};\quad publish $e_t$ to the shared registry\;
$\mathcal{O}_t\gets$ tracked clusters that are moving and small \eqref{eq:track}\;
$\alpha_t\gets$ gain from $e_t$ and $\mathrm{TTC}_t$ by \eqref{eq:alpha}\;
\eIf{$\mathcal{O}_t=\emptyset$}{$\mathbf{u}^{\mathrm{safe}}_t\gets\mathrm{clip}_{\mathcal{U}}(\mathbf{u}^{\mathrm{nom}}_t)$\;}{
  keep the $N_{\mathrm{obs}}$ obstacles of smallest clearance; solve the QP \eqref{eq:qp} with gain $\alpha_t$\;
  \If{solver failed \textbf{or} $\delta>\epsilon_\delta$}{$v^{\mathrm{safe}}_t\gets0$ \tcp*{brake, keep the turn}}
}
\Return{$\mathbf{u}^{\mathrm{safe}}_t$}\;
\end{algorithm}

\subsection{Packaging for Nav2: Two Plugins over a Plugin-Free Core}\label{sec:impl}

\sname{} is built on ROS~2 Jazzy and Nav2 \citep{nav2}. The package \code{nav2_se_controller} ships three shared libraries: a plugin-free \code{se_mppi_core} (entrapment, CBF filter, coordinator, tracker, repulsion, gap search), the \code{escape_critic} MPPI critic, and the \code{safe_escape_controller}. Separating the core from the two plugin libraries resolves a \code{class_loader} factory conflict that otherwise prevents the critic and the controller from loading together. Two integration details are easy to get wrong. First, Nav2's \code{CriticManager} resolves a critics-list entry \code{NAME} to the class \code{mppi::critics::NAME}; a critic in any other namespace passes a direct pluginlib unit test yet never loads from the critics list, so \code{EscapeCritic} lives in \code{mppi::critics}. Second, the Jazzy critic data are \code{xtensor} tensors, not the Eigen types of Nav2's development branch, so the critic is written against the installed headers. The QP uses \code{osqp-eigen} \citep{osqp}, and the build is reproducible through RoboStack (conda-forge) independently of the host OS. The modules this paper relies on (entrapment detector, CBF filter, coordinator, tracker, gap search, repulsion, path progress and plugin loading) carry 44 gtest cases; with the follow-up prediction and multi-robot modules and the live-injection regression of \cref{sec:live}, the package holds 82 \code{TEST} cases in 13 files.
